\documentclass[12pt]{article}
% Shared print layout. The web's provisional design and MathJax are independent.
\usepackage[utf8]{inputenc}
\usepackage[T1]{fontenc}
\usepackage[spanish,es-nodecimaldot]{babel}
\usepackage{lmodern}
\usepackage{amsmath,amssymb}
\usepackage{graphicx}
\usepackage{geometry}
\geometry{a4paper,left=20mm,right=20mm,top=25mm,bottom=22mm,headheight=16pt,headsep=8mm}
\usepackage{fancyhdr,lastpage}
\usepackage{enumitem}
\usepackage{xcolor}
\usepackage{pgfplots}
\pgfplotsset{compat=1.18}
\usepgfplotslibrary{fillbetween}
\usepackage[hidelinks,pdfencoding=auto]{hyperref}
\hypersetup{pdflang=es-ES,pdfauthor={},pdfsubject={Matemáticas Aplicadas a las Ciencias Sociales II}}
\pagestyle{fancy}
\fancyhf{}
\fancyhead[L]{\small Matemáticas Aplicadas a las Ciencias Sociales II}
\fancyhead[R]{\small 2.\textsuperscript{o} Bachillerato}
\fancyfoot[L]{\footnotesize Programación lineal II | Edición revisada 1.0}
\fancyfoot[R]{\footnotesize Página \thepage\ de \pageref{LastPage}}
\renewcommand{\headrulewidth}{0.4pt}
\setlength{\parindent}{0pt}
\setlength{\parskip}{0.5em}
\setlist{itemsep=0.3em,topsep=0.3em}
\newcommand{\apartado}[1]{\subsection*{#1}}
\pgfplotsset{
  revision/.style={
    width=0.94\linewidth,height=7.5cm,
    axis lines=middle,grid=major,grid style={gray!18},
    tick label style={font=\small},label style={font=\small},
    xlabel={$x$},ylabel={$y$},samples=160,
    legend style={at={(0.5,-0.16)},anchor=north,draw=none,font=\small},
    clip=true
  }
}

\hypersetup{pdftitle={Programación lineal II - Examen revisado}}
\begin{document}
\begin{minipage}{0.17\linewidth}
\includegraphics[width=0.88\linewidth]{logo-al-andalus.png}
\end{minipage}\hfill
\begin{minipage}{0.79\linewidth}
{\Large\bfseries Programación lineal II}\\[0.4em]
Duración: \textbf{55 minutos}.\\
Puntuación máxima: \textbf{10 puntos}.
\end{minipage}

Apellidos y nombre: \hrulefill

Fecha: \hrulefill

Responde de forma razonada y numera las páginas de tu respuesta. Se permite el uso de calculadora. \textbf{Los tres ejercicios son obligatorios.}
\section*{Ejercicio 1 \hfill 4 puntos}
\textit{Tiempo orientativo: 25 minutos.}

Una empresa de recambios industriales produce baterías de tipos A y B. La producción semanal debe ser de al menos 10 baterías en total y el número de baterías B no puede superar en más de 10 unidades al de baterías A.

Cada batería A tiene unos gastos de producción de 150 euros y cada batería B, de 100 euros. Se dispone de un máximo de 6\,000 euros semanales para este fin.

La empresa vende todo lo que produce; cada batería A genera un beneficio de 130 euros y cada batería B, de 140 euros. Determina:
\begin{enumerate}[label=(\alph*)]
\item Cuántas baterías de cada tipo debe producir a la semana para que el beneficio total sea máximo.
\item Cuál es dicho beneficio máximo.
\end{enumerate}
\textit{Los ejercicios 2 y 3 figuran en la página siguiente.}
\clearpage
\section*{Ejercicio 2 \hfill 3 puntos}
\textit{Tiempo orientativo: 15 minutos.}

La Agencia Espacial Europea contará con un presupuesto de 2,4 millones de euros para financiar misiones sobre Observación de la Tierra y programas de Transporte Espacial. Cada misión supone una inversión de 200\,000 euros y cada programa, de 100\,000 euros.

En la decisión final deben \textbf{superarse los 2 millones de euros} de inversión. El número de misiones debe ser al menos 4, pero no más de la mitad del número de programas.

\textbf{Formula, sin resolver}, el problema que permite obtener cuántas misiones y programas se deben llevar a cabo para maximizar $F(x,y)=0{,}6x+0{,}4y$, siendo $x$ el número de misiones e $y$ el de programas.
\section*{Ejercicio 3 \hfill 3 puntos}
\textit{Tiempo orientativo: 15 minutos.}

Representa el recinto limitado por las restricciones, calculando sus vértices:
\[
x+4y\geq5,\qquad x+2y\geq4,\qquad7x+5y\leq35,\qquad x\geq0.
\]
\textquestiondown{}En qué punto de la región la función $F(x,y)=2x+y$ alcanza el \textbf{mínimo} y cuál es dicho valor?
\end{document}
